%% ============================================================
%%  SECCIÓN 10 — Gráficas
%%  Responsable: Minilux
%% ============================================================

\section{Gráficas}
\label{sec:graficas}

\begin{figure}[H]
    \centering
    \includegraphics[width=0.85\linewidth]{grafica_gamma_m1}
    \caption{Coeficiente de tensión superficial $\gamma_i$ para cada
    medición del primer método. La línea roja punteada indica el promedio
    $\overline{\gamma}_1 = \qty{39.73}{\milli\newton\per\meter}$ y la línea
    negra el valor bibliográfico
    $\gamma_{\text{ref}} = \qty{72.8}{\milli\newton\per\meter}$
    \cite{serway}. Figura generada con Python \cite{script_graficas}.}
    \label{fig:grafica_gamma_m1}
\end{figure}

La Fig.~\ref{fig:grafica_gamma_m1} evidencia una dispersión estadística mínima de las mediciones en torno a su valor medio ($\overline{\gamma}_1$). Sin embargo, existe una marcada desviación sistemática por debajo del valor referencial ($\gamma_{\text{ref}}$), cuyas posibles causas experimentales se discuten en la Sección~\ref{sec:conclusiones}.